% !TEX root = main.tex
\section{Other Formulas}
\label{sec:other}

A gap transfers along a kind of embedding of webs. This section uses it
to carry the gap of Theorem~\ref{thm:pentagon} to every formula with a
negative occurrence of a suitable tensor.

\subsection{Coordinate Retracts}
\label{sec:other-retract}

\begin{definition}[coordinate retract\leanref{retract}]
\label{def:retract}
For given formulas $D$ and $A$, the \emph{coordinate retracts} of $D$
into $A$ are the injective maps $j : |D| \to |A|$ such that the
pushforward $j_{*}$, defined by $(j_{*}x)(j(d)) = x_d$ and value $0$
outside the image of $j$, maps $P(D)$ into $P(A)$, and the pullback
$z \mapsto z \circ j$ maps $P(A)$ into $P(D)$. We say that $D$ is a
coordinate retract of $A$ if such a map exists.
\end{definition}

\noindent
The pullback of $j_{*}x$ is $x$, and this allows us to prove that gaps
transfer along coordinate retracts.

\begin{theorem}[transfer\leanref{retract-gap}\leanref{retract-ops}]
\label{thm:retract}
(i) If $D$ is a coordinate retract of $A$ and $D$ has a gap, then $A$ has
a gap.
(ii) Coordinate retracts compose. If $D$ is a coordinate retract of $A$,
then $D^{\perp}$ is a coordinate retract of $A^{\perp}$. The formula $D$
is a coordinate retract of $D \mathbin{\&} B$, of $B \mathbin{\&} D$,
and, if the web of $B$ is nonempty, of $D \otimes B$ and of $B \otimes D$.
\end{theorem}

\noindent
For (i), take $x \in P(D)$ outside the deterministic hull of $D$. If $j_{*}x$ were a
mixture of certain states of $A$, pulling the mixture back along $j$
would write $x = (j_{*}x) \circ j$ as a mixture of certain states of $D$.
For negation, the pairing satisfies
$\langle j_{*}u, z \rangle = \langle u, z \circ j \rangle$, which exchanges
the two conditions of Definition~\ref{def:retract}. For $D \otimes B$,
the map is $j(d) = (d, b)$ for a point $b$ of $B$, so
$j_{*}x = x \otimes \delta_b$, where $\delta_b$ is the indicator of $b$. The pullback condition uses that every element of $P(B)$ is at most
$1$ at $b$ (Theorem~\ref{thm:mall}).

\subsection{Negative Occurrences and Four-Cycles}
\label{sec:other-occ}

With positive and negative occurrences as defined in
Section~\ref{sec:exact-beyond}, we prove the following.

\begin{theorem}[the gap at other formulas\leanref{gap-of-occ}]
\label{thm:occ-gap}
Let $n, m, n', m' \ge 2$, and let $A$ be a formula all of whose data types
are nonempty. If $(\mathbf{n} \multimap \mathbf{m}) \otimes
(\mathbf{n'} \multimap \mathbf{m'})$ occurs negatively in $A$, then $A$
has a gap.
\end{theorem}

\noindent
A positive occurrence of $D$ gives a coordinate retract of $D$ into $A$,
and a negative one a coordinate retract of $D^{\perp}$, by
Theorem~\ref{thm:retract}(ii) and induction on $A$, using that $D$ is a
coordinate retract of $D^{\perp\perp}$. Theorem~\ref{thm:retract}(i) and
Theorem~\ref{thm:pentagon} then give the claim. Examples are $T$ itself,
the type $((\mathbf{n} \multimap \mathbf{m}) \otimes (\mathbf{n'}
\multimap \mathbf{m'})) \multimap B$ of a program that receives two
functions and returns a value of type $B$, for every $B$ with nonempty
data types (Lean checks the instance $B = \mathbf{2}$). A positive occurrence of the core
does not suffice: the core itself has no gap
(Theorem~\ref{thm:second-exact}(i)).

The occurrence condition of Theorem~\ref{thm:occ-gap} refers to one
syntactic subformula. We also prove a version that depends only on the
coherence graphs of the factors.

\begin{theorem}[four-cycles\leanref{retract-tens}\leanref{c4-retract}\leanref{gap-of-c4}]
\label{thm:c4}
If $X$ and $Y$ both have an induced $4$-cycle, then
$(X \otimes Y)^{\perp}$ has a gap, and so does every formula over
nonempty data types in which $X \otimes Y$ occurs negatively.
\end{theorem}

\noindent
The four points of an induced $4$-cycle of $X$ form a coordinate retract of
$\mathbf{2} \multimap \mathbf{2}$ into $X$, with input $0$ sent to $a$ and
$c$ and input $1$ to $b$ and $d$. The pushforward lands in $P(X)$ because
the indicators of the coherent pairs of the cycle lie in $P(X)$ (Theorem~\ref{thm:mall}). The pullback lands in
$P(\mathbf{2} \multimap \mathbf{2})$ because the incoherent pairs have
total weight at most $1$ (Theorem~\ref{thm:mall}). Coordinate retracts
pass through $\otimes$, so Theorems~\ref{thm:pentagon}
and~\ref{thm:retract} give the claim. The coherence graph of
$\mathbf{n} \multimap \mathbf{m}$ has an induced $4$-cycle iff
$n, m \ge 2$, so Theorem~\ref{thm:occ-gap} is the case
$X = \mathbf{n} \multimap \mathbf{m}$ and
$Y = \mathbf{n'} \multimap \mathbf{m'}$. The curried type $(\mathbf{2} \multimap
\mathbf{2}) \multimap ((\mathbf{2} \multimap \mathbf{2}) \multimap
\mathbf{2})$, which Theorem~\ref{thm:occ-gap} does not cover, is an
instance\leanref{gap-curried}.
